% Recommended insertion for Paper A (paper/jga/manuscript.tex), about 1.5 pages.
% Suggested place: a new subsection at the end of Section~\ref{sec:heat} (after Section~\ref{sec:basis}),
% or as the first part of Section~\ref{sec:open}. Uses Paper A's macros and bibliography keys only.
% Full statements and proofs: theory/varcurv/varcurv.tex; checks: theory/varcurv/*.py.

\subsection{Variable curvature}\label{sec:varcurv}

Heat invariants stay local when the curvature varies: $c_j=S_{j-2}+\sum_ia_{j-2}(p_i)$, where
$S_l=\frac1{4\pi}\int_\Orb u_{l+1}\,dA$ integrates a universal polynomial in $K$ and its derivatives and the
cone term $a_l(p)=\frac1m\sum_{j=1}^{m-1}b_l(\Phi^j)$ depends only on the germ of the metric at $p$
\cite{donnelly1976}, \cite[4.1, Thm~4.8]{dggw2008}. Two things change, and they have different effects.

\emph{The cone term.} Write $C_\phi^2=4\sin^2(\phi/2)$ and $\Pi_i(m)=\frac1m\sum_{j=1}^{m-1}C_{2\pi j/m}^{-2i}$, so
that $m\Pi_i(m)$ is an even polynomial of degree $2i$ vanishing at $m=1$.

\begin{proposition}[Cone terms in variable curvature]\label{prop:varcone}
For a cone point $p$ of order $m\ge\max(2,l-1)$,
$a_l(p)=\sum_{i=1}^{l+1}\beta_{l,i}(p)\Pi_i(m)$, with universal polynomials $\beta_{l,i}$ of weight $2l$ in $K$
and its covariant derivatives at $p$. The coefficient of the new power $m^{2l+1}$ is
$\frac{|B_{2l+2}|}{(2l+2)!}\beta_{l,l+1}(p)$, where
\[
\beta_{l,l+1}=4^l\,l!\,\bigl[v^{2l}\bigr](f^{-1})'(v)
=\tfrac{(2l)!}{l!}K^l+\dots+\tfrac{2}{(l-1)!}(-\Delta)^{l-1}K ,
\]
$f(r)$ the length of the distance circle of radius $r$ divided by $2\pi$ (for a rotationally symmetric germ;
in general averaged over directions). Thus $\beta_{2,3}=12K^2-2\Delta K$, $\beta_{3,4}=120K^3-42K\Delta K+\Delta^2K$,
and
\[
a_3(p)=\tfrac{(m^2-1)(m^2+3)(3m^4+2m^2+19)}{30240\,m}K^3
-\tfrac{(m^2-1)(7m^6+47m^4+173m^2+733)}{201600\,m}K\Delta K
+\tfrac{(m^2-1)(m^2+11)(3m^4+10m^2+227)}{3628800\,m}\Delta^2K .
\]
\end{proposition}

So one new odd power of the order still enters per order in $t$, but from order $t^2$ on its coefficient
is not a multiple of $K(p)^l$: it contains $\Delta^{l-1}K(p)$, already visible as the term
$-\frac1{15120}m^5\Delta K$ in Schueth's $t^2$ coefficient \cite[Thm~4.1]{schueth2019}. The top coefficient comes
from following the substitution in Schueth's proof of \cite[Thm~3.7]{schueth2019} to all orders: only the
leading Gaussian and the area element reach the highest pole in $C_\phi$. The $t^3$ formula is our own
computation, checked against \cite[Thm~3.7, Thm~4.1]{schueth2019} at lower order and against \eqref{eq:bl}
at constant curvature.

\emph{The smooth term.} At constant curvature $S_l=\alpha_l\Area/4\pi$ is fixed by $c_1$. In general $S_l$
is a new real number at every order, and it decides everything.

\begin{proposition}[Extension and obstruction]\label{prop:varcurv}
Let $\Orb,\Orb'$ be closed orientable cone-point orbifolds.
\begin{enumerate}[label=(\roman*)]
\item Let $\kappa\ne0$. If at every cone point of both, $K=\kappa$ and $\nabla^jK=0$ for $1\le j\le2L-6$, and
$S_l(\Orb)=S_l(\Orb')$ for $l\le L-2$, then $H_L(\Orb)=H_L(\Orb')$ if and only if \eqref{eq:red1} holds. So
Lemma~\ref{lem:sigdata}, and everything deduced from it, holds in this class with $\Area$ read as
$-2\pi\chi$. For $\kappa=0$ it fails.
\item For any $L$ and any $\kappa\in\mathbb R$, the signatures $(g;m)$ and $(g';m')$ carry metrics of
curvature $\kappa$ near their cone points with the same $H_L$ if and only if
$\frac{2-2g}6+\sum_i\frac{(m_i-1)^2}{12m_i}=\frac{2-2g'}6+\sum_j\frac{(m'_j-1)^2}{12m'_j}$; for instance
$S^2(2^8)$ and $S^2(3,3,3)$.
\item If $\sum_i(1-\frac1{m_i})=\sum_j(1-\frac1{m'_j})>1$ and $\sum_i(m_i-\frac1{m_i})=\sum_j(m'_j-\frac1{m'_j})$,
then for every $g$ the signatures $(g;m)$ and $(g;m')$ carry metrics, flat near the cone points, whose
heat invariants agree to all orders. This includes $\Orb(2,8,8)$ and $\Orb(3,3,12)$
(Theorem~\ref{thm:intro-rigid}(ii)).
\end{enumerate}
\end{proposition}

\begin{proof}[Sketch]
(i) By locality and weight, $a_l(p)$, $l\le L-2$, sees only $\nabla^jK(p)$ with $j\le2l-2$, so it equals
$\kappa^lp_l(m)/m$ \cite{ucar2017}; $S_0=\chi/6$, and Lemma~\ref{lem:triangular} applies. For $\kappa=0$ the cone
terms vanish beyond $t^0$, and (iii) gives counterexamples. (ii) $c_2=\chi/6+\sum_i(m_i^2-1)/(12m_i)$ is
topological \cite[(5.7)]{dggw2008}. Conversely, small conformal bumps $e^{2\lambda\psi}$ placed in a flat disc
change $S_l$ by $\lambda^2F_{l+1}\frac1{2\pi}\int\psi\,\Delta^{l+1}\psi+O(\lambda^3)$ with
$F_n=(-1)^nn(n-1)n!/(2n+1)!\ne0$ (second-order Duhamel). Rescaling a bump by $\varepsilon$ multiplies its effect
on $S_l$ by $\varepsilon^{-2l}$, and $L$ bumps at distinct scales give a local diffeomorphism onto
$\mathbb R^L$ (a generalised Vandermonde). Rescaling the whole configuration then reaches any target,
while a flat cylinder of adjustable length fixes the area. (iii) On $\hat{\mathbb C}$ the flat metric
$|\prod_i(z-z_i)^{1/m_i-1}(z-z_0)^{\alpha_0-1}dz|$, $\alpha_0=\sum_i(1-\frac1{m_i})-1$, has the cone points of
orders $m_i$ and an exact flat cone of angle $2\pi\alpha_0$ at $z_0$. Replace a neighbourhood of $z_0$ by
one fixed smooth surface of genus $g$, the same for $m$ and $m'$, containing a flat cylinder of adjustable
length. Flat regions and flat cone points contribute nothing to $c_j$, $j\ge3$; $c_2$ agrees by hypothesis
and $c_1$ by the cylinder.
\end{proof}

Constant curvature is therefore the setting in which the counts of this paper are about the cone points:
there the smooth part of every invariant is a known multiple of the area, and the triangular structure of
Section~\ref{sec:basis} is all that remains. With variable curvature the cone orders enter in exactly the
same way, as Proposition~\ref{prop:varcone} and (i) show, but they are entangled with the smooth parts, which carry no
structure. Short of knowing those, heat invariants of any finite order hear only $c_2$, and for pairs such as
$\Orb(2,8,8)$, $\Orb(3,3,12)$ they hear no more at any order. Whether $c_2$ alone always allows agreement to all
orders, and whether such pairs can be isospectral, is open.
